% Main-paper replacement for the supplied P1 section.
% Requires amsmath and the bibliography entries used in the original text.
% Include p1-appendix.tex after the paper's existing \appendix command.

\subsection{Dependency-Driven Pipeline (P1)}
\label{sec:p1}
P1 discovers parallelism in irregular graph programs by analyzing the
data, memory, and control dependences of lowered LLVM IR. Its central
abstraction is a phase- and scope-aware effect representation: it records
which graph locations a computation accesses and where that computation
occurs within a traversal. P1 uses this representation to establish
spatial independence between work partitions or temporal independence
between rounds or phase components, then constructs an execution plan
that preserves the remaining dependences.

\paragraph{Graph-loop effects.}
A canonical graph loop traverses a frontier or the full vertex domain
$F_t$. For each source $u$, it executes a preamble $P(u)$, a body $B(u,v)$
for each neighbor $v\in N(u)$, and an epilogue $Q(u)$:
\begin{equation}
E_t = \mathrm{Seq}_{u\in F_t}\Bigl(
P(u)\triangleright
\bigl(\mathrm{Seq}_{v\in N(u)} B(u,v)\bigr)
\triangleright Q(u)\Bigr).
\label{eq:effect-expr}
\end{equation}
Here $\mathrm{Seq}$ retains the serial enumeration order and
$\triangleright$ records phase order within a source iteration. Effects
within a phase form a multiset, written using structural composition
$\parallel$; this notation does not assert that those effects can run
concurrently. Retaining order follows the motivation of sequential
effect systems~\cite{10.1145/2429069.2429074} and effect
quantales~\cite{gordon2017genericapproachflowsensitivepolymorphic,10.1145/3450272}.

Each effect atom records an operation kind, canonical memory base,
index region, and occurrence scope. The kinds distinguish reads, writes,
old-value-dependent updates, recognized algebraic updates, claims, and
frontier activations. Regions distinguish source-derived ($U$),
destination-derived ($V$), data-derived ($D$), and global ($G$) indices;
scopes distinguish once-only, per-source, and per-pair operations.
Fixed-point provenance analysis supplies the regions, and subsequent
analysis annotates reads with temporal information and a live or snapshot
value source. Region tags describe access provenance; they do not by
themselves establish disjointness. Unsupported operations prevent the
parallel rewrite. Appendix~\ref{app:p1-effects} gives the recognition and
temporal-classification rules.

\paragraph{Establishing independence.}
P1 constructs conservative RAW, WAR, and WAW conflict relations from
memory-access footprints, including shared runtime resources. The
spatial analysis considers work instances and their proposed ownership
partitions. The temporal analysis considers rounds and components of
the source phases. These are distinct choices of execution unit, each
with its own dependence relation.

For either axis $A\in\{S,T\}$, let $X$ be the relevant execution units,
$\mathsf{Dep}_A$ their remaining dependence relation, and $\mathcal P$ a
proposed non-trivial partition. The independence criterion is
\begin{equation}
\mathsf{Dep}_A\cap\mathsf{Cross}(\mathcal P)=\varnothing,
\qquad
\mathsf{Cross}(\mathcal P)=
\bigcup_{P_i\ne P_j}(P_i\times P_j).
\label{eq:p1-independence}
\end{equation}
Thus, different parts may execute concurrently when no remaining
dependence crosses the partition. For example, assigning all neighbor
work for a source to one owner keeps that source's $A[u]$ updates
together; writes through $A[v]$ can still conflict across source owners
and require a separate justification. Source ownership alone is therefore
insufficient to parallelize a loop that also mutates destination state.

P1 may discharge conflicts using applicable transformations, including
snapshotting, staged claims, and privatized updates or reductions.
Each transformation has two obligations: it must preserve the original
loop's observable behavior under its applicability conditions, and its
resulting effects must justify the proposed concurrency. A snapshot must
supply the value required by the original computation; a recognized
update operator alone does not establish that regrouping updates is
valid. Private updates additionally require a valid combination operation
and no unsupported observation of uncombined state. Conflict edges are
removed only when the corresponding obligation has been discharged.

Disconnected components of the remaining interference graph provide
candidate independent groups. Independence on one axis permits concurrency
on that axis while retaining the ordering required on the other; it does
not license arbitrary reordering of the whole loop. The correctness
argument is conditional on sound effect footprints, behavior-preserving
transformations, and a realization that preserves all remaining ordering
and occurrence-scope requirements. Appendix~\ref{app:p1-independence}
makes these conditions explicit.

\paragraph{Constructing an execution plan.}
Given the original loop $L$, its effects $E_t$, and the analysis results
and justified transformations $\mathcal A$, the compiler constructs
\[
\mathcal R(L,E_t,\mathcal A)=\Pi_t.
\]
The plan identifies executable work units, their domains and state, and
their ordering constraints. It is emitted only when the runtime supports
the required realization; otherwise, P1 retains serial execution. Effect
atoms guide this construction, while the original IR supplies the
computation and control flow.

Spatial execution groups work by source or destination ownership using
Graptor's CleanCut layout~\cite{10.1145/3392717.3392753}. Source ownership
keeps a source's preamble, neighbor work, and epilogue together.
Destination ownership can split a source's neighbors across partitions;
where supported, separate source-coverage passes preserve per-source
operations, including for sources with no neighbors. Temporal execution
runs independent phase components concurrently within a source iteration,
or independent groups of rounds concurrently, retaining serial order
within each component or group.

When both round independence and spatial partitioning are established,
P1 combines them. For each round $t$, setup $S_t$ precedes every partition
batch $B_{t,j}$, and finalization $F_t$ follows all batches:
\[
S_t\longrightarrow B_{t,j}\longrightarrow F_t.
\]
Independent rounds use separate mutable auxiliary state and may overlap.
The implemented combination supports round parallelism with spatial
partitioning; source-phase parallelism uses its separate realization.

The runtime instantiates these work units as tasks and schedules them
subject to the plan's precedence edges. Task nodes may contain many
effect occurrences; the execution graph is distinct from the interference
graph used to justify it. CPU and GPU backends consume the same analysis
results, with additional device-capability checks for GPU execution.
Currently, each GPU temporal component executes sequentially in one
thread, exposing temporal concurrency between kernels. Unsupported
offload uses a supported CPU realization. Appendix~\ref{app:p1-runtime}
describes callback construction, resource lifetimes, and backend scheduling.
